\documentclass[11pt]{article}
\usepackage[margin=2.5cm]{geometry}
\usepackage{booktabs}
\usepackage{array}

\title{Progress report\\ \large License plate detection and recognition on CCPD2019}
\date{Status as of 17 October 2026, for the 20 October ``second approach'' milestone}

\begin{document}
\maketitle

\section{Milestone status}

The second approach is working: a plate-specific CRNN/CTC recognizer, trained
deterministically on the 3{,}500 original-scale training crops, now
outperforms the pretrained EasyOCR baseline under both recognition protocols.
Detection remains the strong stage; recognition was the weak spot and is now
substantially improved.

\begin{table}[h]
\centering
\begin{tabular}{@{}lll@{}}
\toprule
Stage & Model & Status \\
\midrule
Dataset & CCPD 5{,}000-image subset, 15{,}000 scale-variant images, leak-free splits & complete \\
Detection (E1 baseline) & fine-tuned YOLO11s, frozen & complete \\
Recognition baseline (E2) & pretrained EasyOCR, frozen & complete \\
Second approach (E2/E2b) & CRNN + BiLSTM + CTC, trained on CCPD chars & complete \\
Preprocessing ablation (E3) & none / gray+contrast / Canny & complete \\
Size and tilt sensitivity (E4/E5) & per-bucket and per-bin metrics & complete \\
Confidence threshold (E6) & accuracy--coverage sweep, validation-selected tau & complete \\
Detector fine-tuning comparison (E1) & pretrained vs fine-tuned YOLO11 & remaining \\
\bottomrule
\end{tabular}
\end{table}

\section{Test results}

\subsection{Detection (frozen YOLO11s)}

\begin{table}[h]
\centering
\begin{tabular}{@{}lrrrr@{}}
\toprule
Split & mAP@0.5 & Precision & Recall & F1 \\
\midrule
Validation & 0.9245 & 0.9402 & 0.8986 & 0.9189 \\
Test & 0.9073 & 0.9408 & 0.8938 & 0.9167 \\
\bottomrule
\end{tabular}
\end{table}

\subsection{Recognition (EasyOCR vs CRNN)}

Both engines are evaluated through the same harness: mean per-character
confidence, tau selected on validation by accepted exact accuracy + coverage
$-1$ (both select mode \texttt{none}, tau 0.50), test reported once per
protocol.

\begin{table}[h]
\centering
\begin{tabular}{@{}llrrrrrr@{}}
\toprule
Protocol & Engine & Coverage & Char.\ accuracy & Exact plates & Precision & Recall & F1 \\
\midrule
ocr\_only & EasyOCR & 0.7300 & 0.3474 & 0 / 800 & 0.0377 & 0.0924 & 0.0535 \\
ocr\_only & CRNN & 1.0000 & 0.7538 & 269 / 800 & 0.5500 & 1.0000 & 0.7097 \\
end\_to\_end & EasyOCR & 0.7025 & 0.3240 & 0 / 800 & 0.0267 & 0.0593 & 0.0368 \\
end\_to\_end & CRNN & 0.9762 & 0.6561 & 133 / 800 & 0.3624 & 0.9371 & 0.5226 \\
\bottomrule
\end{tabular}
\end{table}

Per-bucket and per-tilt-bin CRNN numbers, and the validation ablation, are in
\texttt{models/crnn/REPORT.md}.

\section{Observations}

\begin{itemize}
\item CRNN lifts test character accuracy to 0.7538 (ocr\_only) and 0.6561
  (end\_to\_end), with 269 and 133 exact plates where EasyOCR read zero.
\item The end-to-end gap to OCR-only is the detector floor plus box geometry.
\item Large plates degrade sharply (ocr\_only 0.5722 vs 0.8635 for small)
  because the fixed 48-px input downsamples them.
\item The 10--19 degree tilt bin is the weakest recognition bin.
\item Canny preprocessing collapses for the trained model; the selection
  correctly picks \texttt{none} for both protocols.
\end{itemize}

\section{Remaining work}

\begin{itemize}
\item E1: fine-tune YOLO11 on the train split and compare against the frozen
  detector; re-run end-to-end recognition with the better detector.
\item Large-plate input geometry (taller CRNN input) as a recognition
  improvement.
\item Final submission: failure-case analysis, latency/memory logging,
  clean-clone reproducibility check, two-page summary and demo
  (10 November).
\end{itemize}

\end{document}
